\documentclass[12pt]{article}

\usepackage{amsmath, amssymb}
\usepackage{geometry}
\usepackage{booktabs}
\usepackage{setspace}

\geometry{margin=1in}
\setstretch{1.2}

\title{CSE400 -- Fundamentals of Probability in Computing\\
Lecture L15--L18: Joint Random Variables and Joint Distributions}
\author{Dhaval Patel, PhD\\
Associate Professor -- Computer Science and Engineering\\
SEAS, Ahmedabad University}
\date{March 10, 2026}

\begin{document}

\maketitle

\section*{1. Introduction / Scope}

In science and in real life we are often interested in \textbf{two (or more) random variables at the same time}.

Examples given in the lecture include:

\begin{itemize}
\item Height and weight of giraffes
\item IQ and birthweight of children
\item Frequency of exercise and rate of heart disease
\item Air pollution level and respiratory illness rate
\item Number of Facebook friends and age of members
\end{itemize}

Additional application examples discussed in the lecture:

\subsection*{Example 1: Smart Transportation}

Two random variables are defined:

\begin{itemize}
\item $X$ = Vehicle Speed
\item $Y$ = Traffic Density
\end{itemize}

Question posed in the lecture:

\begin{quote}
Can speed be high when traffic density is also high?
\end{quote}

Applications mentioned:

\begin{itemize}
\item Adaptive traffic signals
\item Congestion prediction
\item Autonomous vehicle routing
\end{itemize}

\subsection*{Example 2: Wireless Network Performance}

Two random variables:

\begin{itemize}
\item $X$ = Signal-to-Noise Ratio (SNR)
\item $Y$ = Data Throughput
\end{itemize}

Throughput depends on SNR, but may also be affected by network congestion.

Applications mentioned:

\begin{itemize}
\item 5G scheduling
\item Adaptive modulation
\item Network planning
\end{itemize}

\subsection*{Additional Examples}

\begin{itemize}
\item Weather prediction
\begin{itemize}
\item $X =$ Rainfall
\item $Y =$ Temperature
\end{itemize}

\item Drone delivery
\begin{itemize}
\item $X =$ Wind speed
\item $Y =$ Delivery time
\end{itemize}

\item Rolling two dice where conditions such as $X+Y=7$ or $X>Y$ may be studied.
\end{itemize}

\section*{2. Discrete Case -- Joint PMF}

\subsection*{Setup}

Suppose $X$ and $Y$ are discrete random variables.

\[
X = \{x_1, x_2, \ldots, x_n\}
\]

\[
Y = \{y_1, y_2, \ldots, y_m\}
\]

The ordered pair $(X,Y)$ takes values in the product set

\[
\{(x_1,y_1), (x_1,y_2), \ldots, (x_n,y_m)\}
\]

\subsection*{Definition (Joint PMF)}

The joint probability mass function is defined as

\[
p(x_i,y_j) = P(X=x_i, Y=y_j)
\]

This gives the probability of the \textbf{joint outcome occurring simultaneously}.

\subsection*{Joint PMF Table Representation}

The joint PMF can be organized in a table.

\begin{center}
\begin{tabular}{c c c c c}
\toprule
$X\backslash Y$ & $y_1$ & $y_2$ & $\dots$ & $y_m$ \\
\midrule
$x_1$ & $p(x_1,y_1)$ & $p(x_1,y_2)$ & $\dots$ & $p(x_1,y_m)$ \\
$x_2$ & $p(x_2,y_1)$ & $p(x_2,y_2)$ & $\dots$ & $p(x_2,y_m)$ \\
$\vdots$ & $\vdots$ & $\vdots$ & $\dots$ & $\vdots$ \\
$x_n$ & $p(x_n,y_1)$ & $p(x_n,y_2)$ & $\dots$ & $p(x_n,y_m)$ \\
\bottomrule
\end{tabular}
\end{center}

\subsection*{Key Properties}

\[
0 \le p(x_i,y_j) \le 1
\]

\[
\sum_{i=1}^{n} \sum_{j=1}^{m} p(x_i,y_j) = 1
\]

\section*{3. Example -- Rolling Two Dice}

\subsection*{Experiment}

Roll two fair dice.

Define:

\begin{itemize}
\item $X$ = value on the first die
\item $T$ = total (sum) of both dice
\end{itemize}

Each outcome $(i,j)$ has probability

\[
\frac{1}{36}
\]

A joint probability table of $(X,T)$ is constructed.

\subsection*{Observation}

\begin{itemize}
\item Table entries are either $0$ or $\frac{1}{36}$
\item $X$ and $T$ are \textbf{not independent}
\end{itemize}

\section*{4. Continuous Case -- Joint PDF}

The continuous case is analogous to the discrete case.

\begin{itemize}
\item Discrete values $\rightarrow$ continuous intervals
\item Joint PMF $\rightarrow$ Joint PDF
\item Sums $\rightarrow$ Double integrals
\end{itemize}

If

\[
X \in [a,b], \quad Y \in [c,d]
\]

then the pair $(X,Y)$ takes values in

\[
[a,b] \times [c,d]
\]

\subsection*{Definition (Joint PDF)}

A function $f(x,y)$ is called the \textbf{joint probability density function} of $X$ and $Y$ if

\[
P((X,Y)\text{ lies in a small rectangle of area } dx\,dy) \approx f(x,y)\,dx\,dy
\]

\subsection*{Key Properties}

\[
f(x,y) \ge 0
\]

\[
\int_c^d \int_a^b f(x,y)\,dx\,dy = 1
\]

Note from the lecture:

\begin{quote}
The joint PDF $f(x,y)$ may take values greater than 1. It is a probability density, not a probability.
\end{quote}

\section*{5. Worked Example}

Given

\[
f(x,y)=4xy, \quad 0 \le x \le 1,\quad 0 \le y \le 1
\]

Tasks:

\begin{enumerate}
\item Show that $f(x,y)$ is a valid joint PDF.
\item Visualize event $A=\{X<0.5,\,Y>0.5\}$.
\item Compute $P(A)$.
\end{enumerate}

$(X,Y)$ lies in the unit square

\[
[0,1]^2
\]

\subsection*{Step 1 -- Check Validity (Normalization)}

\[
\int_0^1 \int_0^1 4xy\,dx\,dy
\]

First integrate with respect to $x$:

\[
\int_0^1 4xy\,dx = 4y \int_0^1 x\,dx
\]

\[
=4y \left[\frac{x^2}{2}\right]_0^1
\]

\[
=4y\left(\frac{1}{2}\right)
\]

\[
=2y
\]

Now integrate with respect to $y$:

\[
\int_0^1 2y\,dy
\]

\[
=\left[y^2\right]_0^1
\]

\[
=1
\]

Hence $f(x,y) \ge 0$ and total probability is 1.

Therefore $f(x,y)$ is a \textbf{valid joint PDF}.

\subsection*{Step 2 -- Compute $P(A)$}

Event:

\[
A = \{X < 0.5,\; Y > 0.5\}
\]

\[
P(A)=\int_0^{0.5} \int_{0.5}^1 4xy\,dy\,dx
\]

First integrate with respect to $y$:

\[
\int_{0.5}^{1} 4xy\,dy
\]

\[
=4x\left[\frac{y^2}{2}\right]_{0.5}^{1}
\]

\[
=2x(1-0.25)
\]

\[
=2x(0.75)
\]

\[
=\frac{3}{2}x
\]

Now integrate with respect to $x$:

\[
\int_0^{0.5} \frac{3}{2}x\,dx
\]

\[
=\frac{3}{2}\left[\frac{x^2}{2}\right]_0^{0.5}
\]

\[
=\frac{3}{2}\left(\frac{0.25}{2}\right)
\]

\[
=\frac{3}{16}
\]

Thus

\[
P(A)=\frac{3}{16}
\]

\section*{6. Joint Cumulative Distribution Function (Joint CDF)}

Suppose $X$ and $Y$ are jointly distributed random variables.

Notation used:

\[
X \le x,\; Y \le y
\]

which denotes the event

\[
\{X \le x \text{ and } Y \le y\}
\]

\subsection*{Definition}

The joint cumulative distribution function is

\[
F(x,y) = P(X \le x, Y \le y)
\]

\section*{7. Joint CDF -- Continuous Case}

If $X$ and $Y$ are continuous random variables with joint density $f(x,y)$ over $[a,b]\times[c,d]$,

\[
F(x,y)=\int_c^{y}\int_a^{x} f(u,v)\,du\,dv
\]

\subsection*{Recovering the Joint PDF}

Using partial derivatives:

\[
f(x,y) = \frac{\partial^2}{\partial x\,\partial y} F(x,y)
\]

\section*{8. Joint CDF -- Discrete Case}

If $X$ and $Y$ are discrete random variables with joint PMF $p(x_i,y_j)$,

\[
F(x,y)=\sum_{x_i \le x}\sum_{y_j \le y} p(x_i,y_j)
\]

\subsection*{Interpretation}

Add all probabilities:

\begin{itemize}
\item rows with $x_i \le x$
\item columns with $y_j \le y$
\end{itemize}

\section*{9. Properties of the Joint CDF}

The joint CDF $F(x,y)$ must satisfy the following properties.

\subsection*{Property 1}

\[
F(x,y) \text{ is non-decreasing}
\]

If $x$ or $y$ increases, then $F(x,y)$ must stay constant or increase.

\subsection*{Property 2}

\[
F(x,y)=0
\]

at the \textbf{lower-left corner of the joint range}.

If the lower-left corner is $(-\infty,-\infty)$,

\[
\lim_{(x,y)\to(-\infty,-\infty)} F(x,y)=0
\]

\subsection*{Property 3}

\[
F(x,y)=1
\]

at the \textbf{upper-right corner of the joint range}.

If the upper-right corner is $(\infty,\infty)$,

\[
\lim_{(x,y)\to(\infty,\infty)} F(x,y)=1
\]

\end{document}